\section{Using the Prototype To Create Websites}\label{sec:creation-process-of-websites}
This section focuses on explaining the process of setting up the prototype framework, RouteChasm~\parencite{github_routechasm}, and using it to create
three websites of varying types:

\begin{itemize}
    \item Personal - Chrome browser extension for widget-oriented New Tabs~\parencite{github_startuh}
    \item Corporate - Software Agency advertising their services, projects, and blog articles~\parencite{github_strata}
    \item Content-driven - Administrator curated repository of recipes and cooking related articles~\parencite{github_savora}
\end{itemize}

The websites are testing the following sections:

\begin{itemize}
    \item CMS components
    \item AI generation (Page Structure Generation, Article Generation, Static Page Generation, Page Component Generation, Documentation Generation)
    \item Development metrics (Requirement Coverage, Number of Manual Edits, Time to Publish, Responsiveness, and Basic SEO)
\end{itemize}

\subsection{PHP and MySQL}\label{subsec:php-and-mysql}
The recommended way to host the application locally is by installing Wampserver~\parencite{wampserver}.
Download the full version of Wampserver, this will include PHP server, MySQL server, and PHPMyAdmin to manage the MySQL server.
While installing the Wampserver you might need to install Visual C++ Redistributable Packages for your system which can
also be found on the Wampserver website.

After installing and running it, download cacert.pem file, and add absolute location to the file to the curl.cainfo field
in php.ini file.
This will enable the server side AI generation requests.

\subsection{Initial setup}\label{subsec:initial-setup}
For each website the following shell code was executed to create project directory, download the latest version of
the framework, and change the remote origin (GitHub repository).

\begin{verbatim}
    # Download the framework
    git clone https://github.com/CaptSiro/route-chasm.git <project-name>
    cd <project-name>

    # Remove project's git history (recommended through file explorer)
    rm -rf .git
    # Create the project's local repository
    git init

    # Push the framework to the origin
    git remote add origin <project-git>
    git add .
    git commit -m "Initial commit from RouteChasm"
    git branch -M main
    git push -u origin main
\end{verbatim}

Before the project is ready for local testing a few configuration files needs to be updated.
The .htaccess file routes all the requests to a singular file and the framework takes care of the individual request resolution.
By default, this file routes the requests to the file route-chasm/index.php but the real location of the index.php file
depends on where the project is located in relation to the web server's root directory.
Most of the time, the web server's root directory is the parent directory of the project.

\begin{verbatim}
    ...
    RewriteRule (.*) /<project-name>/index.php [QSA,END]
    ...
\end{verbatim}

Now, any incoming requests are handled but the framework lacks environment specific configuration.
By default, the framework tries to get all the configuration from the .env file.
The file is not provided in the framework repository, because it defines sensitive information, for example, root login details,
database user login, and api keys.
The .env file must be created and edited by hand.
The following .env template contains all the fields used by the framework with sample values.

\begin{verbatim}
    VERSION=1.0.2

    PROJECT=<Project-Name>
    PROJECT_LINK=<project-git>
    PROJECT_AUTHOR=<Project-Author>
    PROJECT_AUTHOR_LINK=<project-author-link>

    LANGUAGE=en-US
    DOMAIN_URL=<deploy-url>

    DATABASE_HOST=127.0.0.1
    DATABASE_NAME=routechasm
    DATABASE_USER=root
    DATABASE_PASSWORD=root

    ADMIN_LOGIN_PASSWORD=root

    OPENAI_KEY=<api-key>
    OPENAI_MODEL=gpt-4o-mini
\end{verbatim}

Login to the database manager (PHPMyAdmin recommended and preinstalled in Wampserver) and create new database.
Locate the src/core/sql/init.sql file and run it in the database.
This will create empty RouteChasm database.
To load in the database with data run the src/core/sql/export.sql file.
Each project includes its own export.sql file.

The framework is ready to be used.

\subsection{Implementation Specifics of Corporate Website Scenario}\label{subsec:implementation-specifics-of-corporate-website-scenario}
The Home directory provided by the framework does not include JavaScript file,
but the project homepage requires extra JavaScript functionality and thus home.js is created
and the functionality is implemented.
The template is edited to import the script file using SideLoader's Javascript::import() method and the homepage layout
is designed and implemented.



\subsection{Common Edits}\label{subsec:common-edits}
This section is going to address a few common changes, that can be applied to the framework.

\subsubsection*{Adding Meta Description to Homepage}
All pages created through administration dashboard include textarea to fill in meta description.
The homepage component is defined in src/components/Home/Home.php.
In the constructor the HtmlHead is created, and it includes function to add meta fields.
Adding description editable through web settings can be added this way:

\begin{verbatim}
    parent::__construct(
        new ContextAwareWebPage(
            // Save the HtmlHead to variable
            head: $head = new HtmlHead(...)
        )
    );

    $description = Setting::fromName(
        'project:homepage-description',
        // Create the setting if it is not defined
        true,
        // Default description
        '',
        // Makes the setting editable in admin dashboard
        [PROPERTY_EDITABLE => true]
    );

    $head->addMeta('description', $description->toString());
\end{verbatim}

\subsubsection*{Editing Homepage Button Behavior}
Locate the src/components/Home/Home.phtml template and edit the HTML of the page.
To add JavaScript functionality locate home.js in the same directory.
If the script file is created it needs to be created by hand and this line of code added to the template:

\begin{verbatim}
    // getRosource function locates the file
    // relative to the class definition of the $this object
    Javascript::import($this->getResource('home.js'));
\end{verbatim}

Edit the script to the needs and reload the page.
The changes should take place instantaneously.

\subsubsection*{Uploading Files}
Files can be uploaded via the administration dashboard.
The files are limited to pages created in dashboard and can not be used natively in homepage template.
Navigate to the Admin/File System/Files page and drag-and-drop your files on to the grid.
Drag-and-dropping is the only way to upload files to the system.

\subsubsection*{Changing Homepage Image}
Add images to the public/images directory, locate the src/components/Home/Home.phtml template, and create or edit the
desired HTML image element.
The images are accessible on url: <project-domain>/public/images/<image-file>

\subsubsection*{Creating Editor User}
Navigate to the System/Groups and create new group called Editors and fill out the privileges mapping.
Each checkbox corresponds to Resource and Action performed on the resource.
If the checkbox is checked it means that the group has privilege to perform the action on the resource.

Navigate to the System/Users and create new user.
When creating a new user, the Tag needs to be unique for the website.
Add the user into the Admin group which gives them the access to log into the administration dashboard
and the newly created Editor group defines all the actions the user can perform in the dashboard.

